\footnotesize
\begin{longtable}{>{\raggedright\arraybackslash}p{0.12\textwidth}>{\raggedright\arraybackslash}p{0.16\textwidth}>{\raggedright\arraybackslash}p{0.30\textwidth}>{\raggedright\arraybackslash}p{0.30\textwidth}}
  \caption{List of all datasets in the \celegans{} Connectome Toolbox}\label{tab:dataset-table}\\
  \toprule%
  \textbf{Original publication} & \textbf{Reference/links} & \textbf{Description} & \textbf{Weight} \\
  \midrule%
  \endfirsthead
  \caption[]{(continued)}\\
  \toprule%
  \textbf{Original publication} & \textbf{Reference/links} & \textbf{Description} & \textbf{Weight} \\
  \midrule%
  \endhead
  \midrule
  \multicolumn{4}{r}{\footnotesize\itshape Continued on next page}\\
  \endfoot
  \bottomrule
  \endlastfoot

  \multirow{3}{\linewidth}{\cite{White1986}}  & White et al. 1986 JSH/L4 \newline \href{https://openworm.org/ConnectomeToolbox/WhiteJSH_data}{Matrix} | \href{https://openworm.org/ConnectomeToolbox/Validation#whiteetal1986}{Validation} \newline  & Chemical and electrical connectivity of the JSH L4 worm from White et al. 1986 data, taken from the neurodata.txt file from R. Durbin's thesis 1987. \newline &  Weights represent the count of individual presynaptic active zones observed between two neurons across all the serial EM micrographs. \\
  & White et al. 1986 N2U/adult \newline \href{https://openworm.org/ConnectomeToolbox/WhiteN2U_data}{Matrix} | \href{https://openworm.org/ConnectomeToolbox/Validation#whiteetal1986}{Validation} \newline  & Chemical and electrical connectivity of the N2U (adult hermaphrodite) worm from White et al. 1986 data, taken from the neurodata.txt file from R. Durbin's thesis 1987. \newline &  \emph{(Same as above)} \\
  & White et al. 1986 (whole worm) \newline \href{https://openworm.org/ConnectomeToolbox/White_whole_data}{Matrix} | \href{https://openworm.org/ConnectomeToolbox/Validation#whiteetal1986}{Validation} \newline  & Reanalysis of the White et al. 1986 connectome data by Varshney et al. 2011, but including the pharynx. \newline &  Used Varshney et al. 2011 interpretation of weights: the total number of synaptic contacts from neuron A to neuron B. Contacts are given equal weight regardless of the apparent size of the synaptic apposition. \\
  \midrule%

  \multirow{1}{\linewidth}{\cite{Varshney2011}}  & Varshney et al. 2011 \newline \href{https://openworm.org/ConnectomeToolbox/Varshney_data}{Matrix} | \href{https://openworm.org/ConnectomeToolbox/Validation#varshneyetal2011}{Validation} \newline  & A corrected and extended version of the White et al. 1986 chemical and electrical wiring diagram, incorporating original Mind of a Worm data, Durbin's unpublished reconstructions, new EM imaging of previously unimaged dorsal cord regions, and over 3,000 synapse additions or corrections, particularly in the ventral cord motor neuron connectivity. Excludes pharyngeal neurons, but includes neuromuscular junction connections, all to one BWM (body wall muscle cell). \newline &  Weights are the total number of synaptic contacts from neuron A to neuron B. Contacts are given equal weight regardless of the apparent size of the synaptic apposition. \\
  \midrule%

  \multirow{2}{\linewidth}{\cite{Bentley2016}}  & Bentley et al. 2016  (monoamin.) \newline \href{https://openworm.org/ConnectomeToolbox/Bentley2016_MA_data}{Matrix} | \href{https://openworm.org/ConnectomeToolbox/Validation#bentleyetal2016}{Validation} \newline  & Data on monoaminergic connectivity from Bentley et al. 2016 (i.e. dopaminergic, tyraminergic, octopaminergic \& serotonergic extracellular transmission). \newline &  Adjacency matrices are binary and directed; a weight of 1 between neurons A and B signifies that cell A expresses a biosynthetic enzyme or transporter for the specific monoamine and neuron B expresses a cognate receptor for that monoamine. \\
  & Bentley et al. 2016 Pep \newline \href{https://openworm.org/ConnectomeToolbox/Bentley2016_Pep_data}{Matrix} | \href{https://openworm.org/ConnectomeToolbox/Validation#bentleyetal2016}{Validation} \newline  & Data on peptidergic connectivity from Bentley et al. 2016 (i.e. extracellular synaptic transmission via neuropeptides). \newline &  Adjacency matrices are binary and directed; a weight of 1 between neurons signifies that a ligand-receptor pair with a biologically plausible EC50 is co-expressed across the two neurons. \\
  \midrule%

  \multirow{2}{\linewidth}{\cite{Cook2019}}  & Cook et al. 2019 (herm.) \newline \href{https://openworm.org/ConnectomeToolbox/Cook2019Herm_data}{Matrix} | \href{https://openworm.org/ConnectomeToolbox/Validation#cooketal2019}{Validation} \newline  & Chemical and electrical connectivity of the hermaphrodite from Cook et al. 2019, including connections between neurons, muscles and other cells. \newline &  Weights are the total number of EM serial sections of connectivity, taking into account both the number of synapses and the sizes of synapses. \\
  & Cook et al. 2019 (male) \newline \href{https://openworm.org/ConnectomeToolbox/Cook2019Male_data}{Matrix} | \href{https://openworm.org/ConnectomeToolbox/Validation#cooketal2019}{Validation} \newline  & Chemical and electrical connectivity of the male from Cook et al. 2019, including connections between neurons, muscles and other cells. \newline &  \emph{(Same as above)} \\
  \midrule%

  \multirow{1}{\linewidth}{\cite{Cook2020}}  & Cook et al. 2020 \newline \href{https://openworm.org/ConnectomeToolbox/Cook2020_data}{Matrix} | \href{https://openworm.org/ConnectomeToolbox/Validation#cooketal2020}{Validation} \newline  & Chemical and electrical connectivity of the \celegans{} pharynx, from Cook et al 2020, including connections between pharyngeal neurons, muscles and other cells (epithelial, gland and marginal). \newline &  Weights represent the anatomical strength of a connection, calculated as the total number of serial EM sections across which the connection is observed, summed over all individual synapses between the two cells. For chemical connections, a section is counted if a presynaptic specialization is visible; for electrical connections, a section is counted if a gap junction is visible between the two cells. \\
  \midrule%

  \multirow{1}{\linewidth}{\cite{Brittin2021}}  & Brittin et al. 2021 \newline \href{https://openworm.org/ConnectomeToolbox/Brittin2021_data}{Matrix} | \href{https://openworm.org/ConnectomeToolbox/Validation#brittinetal2021}{Validation} \newline  & Membrane contact information from Brittin et al. 2021. This dataset contains information on the contact area between pairs of cells in the \celegans{} nerve ring, as measured by electron microscopy. The M$^4$ graph as described in the paper is used here. 
 \newline &  Weights represent the contact area between pairs of cells \\
  \midrule%

  \multirow{8}{\linewidth}{\cite{Witvliet2021}}  & Witvliet et al. 2021 1 (L1) \newline \href{https://openworm.org/ConnectomeToolbox/Witvliet1_data}{Matrix} | \href{https://openworm.org/ConnectomeToolbox/Validation#witvlietetal2021}{Validation} \newline  & Chemical and electrical connectivity of from Witvliet et al. 2021, dataset 1 (L1 stage) \newline &  Weights represent the integer number of agreed presynaptic active zones between each pair of cells, counted from electron micrographs and agreed upon by at least two of three independent annotators \\
  & Witvliet et al. 2021 2 (L1) \newline \href{https://openworm.org/ConnectomeToolbox/Witvliet2_data}{Matrix} | \href{https://openworm.org/ConnectomeToolbox/Validation#witvlietetal2021}{Validation} \newline  & Chemical and electrical connectivity of from Witvliet et al. 2021, dataset 2 (L1 stage) \newline &  \emph{(Same as above)} \\
  & Witvliet et al. 2021 3 (L1) \newline \href{https://openworm.org/ConnectomeToolbox/Witvliet3_data}{Matrix} | \href{https://openworm.org/ConnectomeToolbox/Validation#witvlietetal2021}{Validation} \newline  & Chemical and electrical connectivity of from Witvliet et al. 2021, dataset 3 (L1 stage) \newline &  \emph{(Same as above)} \\
  & Witvliet et al. 2021 4 (L1) \newline \href{https://openworm.org/ConnectomeToolbox/Witvliet4_data}{Matrix} | \href{https://openworm.org/ConnectomeToolbox/Validation#witvlietetal2021}{Validation} \newline  & Chemical and electrical connectivity of from Witvliet et al. 2021, dataset 4 (L1 stage) \newline &  \emph{(Same as above)} \\
  & Witvliet et al. 2021 5 (L2) \newline \href{https://openworm.org/ConnectomeToolbox/Witvliet5_data}{Matrix} | \href{https://openworm.org/ConnectomeToolbox/Validation#witvlietetal2021}{Validation} \newline  & Chemical and electrical connectivity of from Witvliet et al. 2021, dataset 5 (L2 stage) \newline &  \emph{(Same as above)} \\
  & Witvliet et al. 2021 6 (L3) \newline \href{https://openworm.org/ConnectomeToolbox/Witvliet6_data}{Matrix} | \href{https://openworm.org/ConnectomeToolbox/Validation#witvlietetal2021}{Validation} \newline  & Chemical and electrical connectivity of from Witvliet et al. 2021, dataset 6 (L3 stage) \newline &  \emph{(Same as above)} \\
  & Witvliet et al. 2021 7 (adult) \newline \href{https://openworm.org/ConnectomeToolbox/Witvliet7_data}{Matrix} | \href{https://openworm.org/ConnectomeToolbox/Validation#witvlietetal2021}{Validation} \newline  & Chemical and electrical connectivity of from Witvliet et al. 2021, dataset 7 (adult stage) \newline &  \emph{(Same as above)} \\
  & Witvliet et al. 2021 8 (adult) \newline \href{https://openworm.org/ConnectomeToolbox/Witvliet8_data}{Matrix} | \href{https://openworm.org/ConnectomeToolbox/Validation#witvlietetal2021}{Validation} \newline  & Chemical and electrical connectivity of from Witvliet et al. 2021, dataset 8 (adult stage) \newline &  \emph{(Same as above)} \\
  \midrule%

  \multirow{1}{\linewidth}{\cite{Randi2023}}  & Randi et al. 2023 \newline \href{https://openworm.org/ConnectomeToolbox/Randi2023_data}{Matrix} | \href{https://openworm.org/ConnectomeToolbox/Validation#randietal2023}{Validation} \newline  & Data on functional connectivity of \celegans{} from WormNeuroAtlas Python package (values extracted with get\_signal\_propagation\_map() \& get\_signal\_propagation\_q(), with max\_q = 0.05) \newline &  Weights represent the mean post-stimulus calcium response amplitude measured in a downstream neuron following two-photon optogenetic stimulation of an upstream neuron \\
  \midrule%

  \multirow{3}{\linewidth}{\cite{RipollSanchez2023}}  & Ripoll Sanchez et al. 2023 (short range) \newline \href{https://openworm.org/ConnectomeToolbox/RipollSanchezShortRange_data}{Matrix} | \href{https://openworm.org/ConnectomeToolbox/Validation#ripollsanchezetal2023}{Validation} \newline  & Dataset of neuropeptidergic connections only between neurons whose processes overlap in the same neuronal process bundle \newline &  Weights indicate the number of different NPP-GPCR pathways connecting a neuron pair in a given direction. \\
  & Ripoll Sanchez et al. 2023 (mid range) \newline \href{https://openworm.org/ConnectomeToolbox/RipollSanchezMidRange_data}{Matrix} | \href{https://openworm.org/ConnectomeToolbox/Validation#ripollsanchezetal2023}{Validation} \newline  & Dataset of neuropeptidergic connections between any neurons whose processes lie within the same broad body region (head, midbody, or tail) \newline &  \emph{(Same as above)} \\
  & Ripoll Sanchez et al. 2023 (long range) \newline \href{https://openworm.org/ConnectomeToolbox/RipollSanchezLongRange_data}{Matrix} | \href{https://openworm.org/ConnectomeToolbox/Validation#ripollsanchezetal2023}{Validation} \newline  & Dataset of neuropeptidergic connections between any neuron pair regardless of anatomical location \newline &  \emph{(Same as above)} \\
  \midrule%

  \multirow{2}{\linewidth}{\cite{Yim2024}}  & Yim et al. 2024 (non norm.) \newline \href{https://openworm.org/ConnectomeToolbox/Yim2024NonNorm_data}{Matrix} | \href{https://openworm.org/ConnectomeToolbox/Validation#yimetal2024}{Validation} \newline  & Reconstruction of the chemical connectome of the dauer, a distinct developmental stage of \celegans{}. \newline &  Weights represent the contact area (nm$^3$) between pairs of cells \\
  & Yim et al. 2024 \newline \href{https://openworm.org/ConnectomeToolbox/Yim2024_data}{Matrix} | \href{https://openworm.org/ConnectomeToolbox/Validation#yimetal2024}{Validation} \newline  & Reconstruction of the chemical connectome of the dauer, a distinct developmental stage of \celegans{}, with the weights normalized to ease comparison to other datasets. \newline &  Weights represent the contact area between pairs of cells, normalized by the standard deviation of connection weights without the top 5th percentile to remove the bias due to the big outliers \\
  \midrule%

  \multirow{2}{\linewidth}{\cite{Wang2024}}  & Wang et al. 2024 (herm.) \newline \href{https://openworm.org/ConnectomeToolbox/Wang2024Herm_data}{Matrix} | \href{https://openworm.org/ConnectomeToolbox/Validation#wangetal2024}{Validation} \newline  & This dataset for the hermaphrodite \celegans{} contains neurotransmitter expression values from: Wang et al. 2024 with basic anatomical connectivity information from Cook et al. 2019, and monoaminergic receptor expression information from Bentley et al. 2016 \newline &  A weight of 1 indicates that the cell expresses the specific neurotransmitter, in addition to an anatomical connection between the pre- and post-synaptic cells (for Glutamate, Acetylcholine, GABA, Betaine), or (for Dopamine, Serotonin, Tyramine, Octopamine) a connection between a monoaminergic receptor-expressing cell and a monoaminergic cell according to Bentley et al. 2016. \\
  & Wang et al. 2024 (male) \newline \href{https://openworm.org/ConnectomeToolbox/Wang2024Male_data}{Matrix} | \href{https://openworm.org/ConnectomeToolbox/Validation#wangetal2024}{Validation} \newline  & This dataset for the male \celegans{} contains neurotransmitter expression values from: Wang et al. 2024 with basic anatomical connectivity information from Cook et al. 2019, and monoaminergic receptor expression information from Bentley et al. 2016 \newline &  \emph{(Same as above)} \\
  \midrule%
\end{longtable}
\normalsize
